% ======================= 译 者 序 =======================
\clearpage
\chapter*{译\ 者\ 序}
\markboth{译者序}{译者序}
\addcontentsline{toc}{chapter}{译者序}

\noindent
二〇二三年四月，莫斯科，列宁山。巴兰尼科夫教授把这批照片摊在桌上，一共三十七张，每张下面压着一张手写的卡片。

他说：请您先看这些字，不要看卡片。

我看了二十分钟，什么也没看出来。照片是黑白的，物体是长条形的，表面刻着字，字很小，笔画匀细。我认出几个，其余不认识。我以为是某种篆书的变体，说了出来。他没有否认，也没有承认。他把其中一张卡片推到我面前。

卡片上写着：「简十四。长 18.4 cm。金属。第 4 号居住环西侧储物间。」

第 4 号居住环。我问他这是什么地方。他说：月球正面，静海西缘，一处编号 \textsc{mare-tranq-07} 的基地。

那是我第一次听说广寒宫不是一座宫殿。

\section*{一、\ 翻译一本没有正文的书}

这本书与《战争的雷霆》不同。那一本有正文，有论证，有结论；这一本没有。

它有三十七片残简，每片下面三段：释文、出处、校记。释文是原件上刻的字，出处是它在基金会的编号与尺寸，校记是辑校者对字词的判断。除此之外，还有一部导论、三万余字的下编、若干附录。

翻译的难处不在术语，在分寸。

释文是用小篆刻的古汉语，一句话常常十来个字，主语与标点都没有，动词与名词的界限也含糊。若照字面译成白话，会丢掉它的简；若照原样保留，中国读者读起来也会像读天书。

我们最后定的办法是这样的：释文一律照原样保留，一个字也不改，连补字符号也照原样；另外在书末加一篇附录，把三十七片逐条译成白话，读者读不通时翻到后面查。

这不是一个好办法。它让读者读一段要翻两回书。我们试过别的方案，比如把白话夹在释文里，或者把释文整体改写成现代汉语，都不行。前者破坏了原件的样式，后者等于替巴兰尼科夫重写了一遍。

附录四就是为此而设。译文只求读得通，不求好看；凡原件没有说的，译文里也不替他补。凡译文与释文有出入，一律以释文为准。

\section*{二、\ 关于那枚圆片}

译者序里要提一件与本书内容无关、但与本书来历有关的事。

一九六三年那批材料里，有一枚直径不足两厘米的圆片，编号 ФВ-1963-46。简三十六记了它：「一物，圆，径不及二寸，重若干。非金非石。击之不鸣。来源不明，暂不归类。」

这枚圆片我们在《战争的雷霆》的译者序里已经提过一次。那一次提它，是因为它的铼锇同位素比值。

这一次提它，是因为另一件事。二〇二三年，我们请人对它作了一次表面微区分析，发现表面的刻痕是一圈极细的同心环，并非装饰，环距不等，最密处在每毫米一百二十条左右。这个密度超出了任何已知的机械刻写方式。

报告交上去以后，基金会的回复只有一句：「该器物不在本次分析范围之内，请勿重复取样。」

我们没有再取样。这一句附在这里，是留一个记录。

\section*{三、\ 体例}

凡例已详，此处只说三条与中译本有关的。

一、释文照原样，一字一格，不作句读。译者所加的断句用顿点，加处在脚注里注明。

二、补字作方框（□），缺字不可知者亦作方框，两者在正文里看不出区别，须看校记。这是原书的做法，我们没有改。

三、书中一切「据《战争的雷霆》」，指《战争的雷霆》的中译本（二〇二一年据全文重译，二〇二四年第四版）。凡引那一本的页次，一律用第四版。

\section*{四、\ 一件没有解决的事}

最后说一件事，说出来对本书不利。

翻译过程中，我们核对了全书所有的数字，凡能互相校验的都校过。有一处始终对不上：简十九记事故二万二千，简二十记实存十一万九千，而完工时的人数按《战争的雷霆》是十五万。

差八千。

巴兰尼科夫在书里承认了这个差额，并给了三种可能的解释，最后说他倾向前者但没有证据。我们照译，没有替他补。

这一处之所以要在这里点出来，是因为它是全书唯一一处作者自己承认不知道的地方。其余的判断，他都很确定；只有这一处，他写了两遍「残简无一句」。

我们认为，这一段应当被中文读者看到。

\begin{flushright}
{\small 李文亚}\quad{\footnotesize\color{pencil} 二〇二四年十月，于北京}
\end{flushright}

% ======================= 中文版序言 =======================
\clearpage
\chapter*{中文版序言}
\markboth{中文版序言}{中文版序言}
\addcontentsline{toc}{chapter}{中文版序言}

\noindent
这本书是为一群我从未见过的人写的。

三十七片残简上的字，是汉字；读它们的人，在中国。而把它们从月面带回地球、编目、封存、六十年无人过问的，是另一群人。我做的事，只是把两头连起来。

连起来之后，我发现一个尴尬的事实：这些字写得很坏。

内容并不坏，坏在写法。它们没有文采，没有修辞，没有一句多余的话。记账的只记账，登记的只登记，连最后一天的场面也只用了二十七个字。若把它们当作文章来读，读三页就会困。

但我建议各位读下去，理由只有一个：这是一个月球上的城在最后两万年里，唯一愿意留给自己的东西。

一个有十一万九千人的城，两万年，只留下三十七片。平均五百四十年一片。这个比例本身，比任何一片的内容都更能说明那座城是什么样子。

\section*{一、\ 关于我的读法}

中国读者大约会问：一个俄国人，凭什么读小篆。

这个问题我答过很多次，答案一直不太好听：凭笨功夫。我从一九六三年开始抄这些字，抄了六十年。抄错的地方，比读对的地方多。

至于方法，我在导论里写了五节，此处不重复。只补一句：这本书里凡是我说有把握的地方，都是因为那一片的字迹清楚、上下文能对上、而且能与其他几片互证；凡是这三条缺一条的，我都写了「残简无一句」。

这句话全书出现二十余次。这话写在纸上，是记录，不是场面话。

\section*{二、\ 关于这本书的命运}

最后说一件与学术无关的事。

这一卷在莫斯科印了两次，第一次是一九八〇年代的内部油印本，印数不足百册；第二次是现在这个排印本。两次之间隔了四十年。

四十年里，这批材料没有丢，原因不好听：它被放在库房最里面，几次清退都没轮到。

我把这件事写进结语，不是发牢骚。我想说的是：一本书能不能留下来，与它写得好不好关系不大。这一点，中国的读者大概比我更清楚。

\begin{flushright}
{\small А.\,В.\,巴兰尼科夫}\quad{\footnotesize\color{pencil} 二〇二四年十一月，于莫斯科}
\end{flushright}
